\section{Background}

The global energy landscape is undergoing a paradigm shift, driven by the imperative to mitigate climate change and the rapid proliferation of renewable energy technologies. Traditional centralized power grids are increasingly being challenged by the integration of Distributed Energy Resources (DERs), such as solar Photovoltaic (PV) systems and Battery Energy Storage Systems (BESS) \cite{ref_iea_2023}. This transition has given rise to the concept of the Microgrid, a localized grouping of electricity generation, energy storage, and loads that normally operates connected to and synchronous with the traditional wide area synchronous grid (macrogrid), but can disconnect and function autonomously as physical island \cite{Zia2018}.

Microgrids offer significant advantages in terms of reliability, efficiency, and sustainability. However, their operation introduces complex control challenges. The primary difficulty arises from the intermittent nature of renewable generation, particularly solar irradiance, which fluctuates stochastically due to weather conditions. This unpredictability creates a mismatch between generation and consumption, necessitating sophisticated Energy Management Systems (EMS) to maintain the power balance and optimize economic dispatch \cite{Olivares2014}.

The emergence of the "Prosumer" (a consumer who also produces energy) has further complicated grid dynamics. In residential settings, prosumers equipped with PV-storage systems face the decision of whether to consume their self-generated energy, store it for later use, or sell it back to the grid. This decision-making process is heavily influenced by dynamic pricing policies, such as Time-of-Use (TOU) tariffs, which are designed to incentivize load shifting away from peak demand periods \cite{Newsham2011}.

Conventionally, EMS optimization has been addressed using rule-based heuristics or mathematical programming techniques like Linear Programming (LP) and Mixed-Integer Linear Programming (MILP) \cite{SARDA2022280}. While optimization-based methods provide theoretically optimal solutions, they typically require accurate forecasts (perfect knowledge) of future load and generation, which are rarely available in real-world scenarios. Furthermore, their computational complexity scales poorly with system size and non-linear battery constraints \cite{Huang2021}. Rule-based methods, conversely, are computationally efficient but inflexible, often failing to adapt to changing grid conditions or pricing structures, leading to suboptimal economic performance.

In recent years, Deep Reinforcement Learning (DRL) has emerged as a promising alternative for real-time microgrid control. Unlike traditional optimization, DRL agents learn control policies directly from interactions with the environment, enabling them to make optimal decisions under uncertainty without requiring an explicit model of the system dynamics \cite{Mocanu2018, Li2019}. Among DRL algorithms, Soft Actor-Critic (SAC) \cite{Haarnoja2018} is particularly well-suited for energy management tasks due to its ability to handle continuous action spaces (e.g., continuous battery charging rates) and its entropy-regularization framework, which encourages robust exploration. This thesis investigates the application of such advanced DRL techniques to bridge the gap between theoretical optimality and practical, real-time applicability in residential microgrids.
